\documentclass[12pt]{article}
\usepackage[utf8]{inputenc}
\usepackage[T1]{fontenc}
\usepackage{hyperref}
\title{Chalikha Language Resources for K-12 Curriculum Development}
\author{Compiled by Claude Sonnet 5 (claude-sonnet-5) for Andr\'as Kornai}
\date{}
\begin{document}
\maketitle

\section{Introduction}
Chalikha (Chali, Chalipkha, Tshali; ISO 639-3 tgf, Glottocode chal1267) is an East Bodish Sino-Tibetan language of Mongar District, eastern Bhutan, related to Bumthangkha and Kurtoepkha. Speaker counts vary widely between sources (roughly 1,400--2,000, with one outlier of 8,200), reflecting different measures.

\section{Primary Educational Resources}
None found.

\section{Linguistic and Reference Materials}
\begin{itemize}
\item Chali language (Wikipedia): \url{https://en.wikipedia.org/wiki/Chali_language}
\item Bhutan Cultural Library text on Chali (Mandala): \url{https://texts.mandala.library.virginia.edu/node/48906}
\item Chalipaikha paper (PDF, Mandala): \url{https://mandala-texts-dev.internal.lib.virginia.edu/sites/mandala-texts-dev.internal.lib.virginia.edu/files/Chalipaikha-ph.pdf}
\item Languages of Bhutan (Wikipedia): \url{https://en.wikipedia.org/wiki/Languages_of_Bhutan}
\end{itemize}

\section{Religious and Community Resources}
\begin{itemize}
\item Joshua Project, Chali in Bhutan: \url{https://joshuaproject.net/index.php/people_groups/18909}
\item Asia Harvest profile: \url{https://people-groups.asiaharvest.org/Bhutan/Chali.pdf}
\end{itemize}

\section{Digital and Audio Resources}
None found.

\section{Curriculum Development Context}
No school materials were found. The Mandala papers may serve as starting points for documentation; any curriculum would need orthography decisions (Tibetan script is listed by Wikipedia) and work with the community.

\end{document}